\chapter{Structure, text and appearance}
\label{ch:text}

\section{Document classes}\label{document-classes}

The class sets the overall type of document.

{\def\LTcaptype{none} % do not increment counter
\begin{longtable}[]{@{}ll@{}}
\toprule\noalign{}
Class & For \\
\midrule\noalign{}
\endhead
\bottomrule\noalign{}
\endlastfoot
\texttt{article} & Papers and short reports \\
\texttt{report} & Long reports with chapters \\
\texttt{book} & Books \\
\texttt{beamer} & Slides \\
\end{longtable}
}

\section{Headings and contents}\label{headings-and-contents}

\begin{lstlisting}[language=latex]
\section{Method}
\subsection{Data}
\subsubsection{Cleaning}
\tableofcontents
\end{lstlisting}

In a complete \texttt{report}, the title page, the contents and each chapter start on new pages:

\textbf{chapters.tex}

\lstinputlisting{../../../teaching/latex/examples/chapters.tex}

Compile it twice, because the contents are only known after the first run.

\begin{figure}
\centering
\includegraphics[width=\linewidth]{img/chapters-title.png}
\caption{Title page}
\end{figure}

\begin{figure}
\centering
\includegraphics[width=\linewidth]{img/chapters-contents.png}
\caption{Contents}
\end{figure}

\begin{figure}
\centering
\includegraphics[width=\linewidth]{img/chapters-chapter.png}
\caption{First chapter}
\end{figure}

Numbers come automatically, and adding a star, as in \texttt{\textbackslash{}section*\{Thanks\}}, leaves a heading unnumbered and out of the contents. Split long documents into files with \texttt{\textbackslash{}input\{chapter1\}}.

\section{Text and fonts}\label{text-and-fonts}

\begin{lstlisting}[language=latex]
\emph{emphasis}  \textbf{bold}  \texttt{code}  \underline{underlined}
{\small smaller text}  {\large larger text}
\end{lstlisting}

Prefer \texttt{\textbackslash{}emph} over forcing italics, because it describes meaning, and it adapts when nested. For dashes, \texttt{-} joins words, \texttt{-\/-} makes ranges such as 10-\/-20, and \texttt{-\/-\/-} is a long dash. Quotes are written as two backticks and two apostrophes.

\section{Fonts}\label{fonts}

With \texttt{pdflatex}, fonts come from packages, for example \texttt{\textbackslash{}usepackage\{helvet\}} with \texttt{\textbackslash{}renewcommand\{\textbackslash{}familydefault\}\{\textbackslash{}sfdefault\}} for a sans-serif body. With \texttt{xelatex}, you can name any installed font:

\begin{lstlisting}[language=latex]
\usepackage{fontspec}
\setmainfont{Source Serif 4}
\end{lstlisting}

\section{Colour}\label{colour}

\begin{lstlisting}[language=latex]
\usepackage{xcolor}
\definecolor{brand}{HTML}{0D6A6E}
\textcolor{brand}{Coloured text}
\end{lstlisting}

Define each colour once, by name, so that you can change the whole document from one line. Don\textquotesingle t rely on colour alone to carry meaning, because some readers can\textquotesingle t see the difference and pages may be printed in black and white.

\section{Lists}\label{lists}

\begin{lstlisting}[language=latex]
\begin{itemize}
  \item A point
\end{itemize}
\begin{enumerate}
  \item A first step
\end{enumerate}
\end{lstlisting}

\section{Page layout}\label{page-layout}

\begin{lstlisting}[language=latex]
\usepackage[a4paper,margin=2.5cm]{geometry}
\usepackage{setspace}
\onehalfspacing
\usepackage{fancyhdr}
\pagestyle{fancy}
\end{lstlisting}

\section{Links}\label{links}

\begin{lstlisting}[language=latex]
\usepackage{hyperref}            % load this near the end of the preamble
\href{https://example.org}{a link}
\url{https://example.org}
\end{lstlisting}

\texttt{hyperref} also makes every cross-reference and contents entry clickable.

